\section{Ambrose's Workshop: Hexapla, Commentaries, First Principles}

\subsection{The patron and the production line}

Around the second decade of the third century Origen acquired the two things
every ancient scholar needed: a converted patron and a staff. Ambrose, a
wealthy man whom Eusebius says Origen won from the Valentinian heresy,
``employed innumerable incentives, not only exhorting him by word, but also
furnishing abundant means'': ``he dictated to more than seven amanuenses, who
relieved each other at appointed times. And he employed no fewer copyists,
besides girls who were skilled in elegant writing,'' all at Ambrose's expense
(\work{Church History} 6.18.1, 6.23.1--2). Origen's output is unintelligible
without this apparatus: it was a publishing operation, and Ambrose kept it
running, and kept its author pressed to produce, for the rest of both their
lives. Jerome, cataloguing the works for Paula near the end of the
fourth century, could compare only Varro, ``because of the countless books
which he wrote,'' with ``our Christian man of brass, or, rather, man of
adamant---Origen, I mean'' (Jerome, \work{Letter} 33.1--3); ancient counts
of the writings ran into the thousands of titles once every homily, scholion,
and letter was numbered.

\subsection{The Hexapla}

The foundation of everything was textual. ``So earnest and assiduous was
Origen's research into the divine words that he learned the Hebrew language,
and procured as his own the original Hebrew Scriptures which were in the
hands of the Jews''---a step without precedent among Christian scholars. He
then collected the rival Greek versions: ``in addition to the well-known
translations of Aquila, Symmachus, and Theodotion, he discovered certain
others which had been concealed from remote times''---one of them, Eusebius
says, found ``in a jar in Jericho in the time of Antoninus, the son of
Severus''; in the Hexapla of the Psalms, ``after the four prominent
translations, he adds not only a fifth, but also a sixth and seventh.'' ``Having collected all of these, he divided them into sections, and
placed them opposite each other, with the Hebrew text itself. He thus left
us the copies of the so-called Hexapla''---the sixfold Bible, which later
writers describe as Hebrew, Hebrew transliterated into Greek letters,
Aquila, Symmachus, the Septuagint, and Theodotion---``He arranged also
separately an edition of Aquila and Symmachus and Theodotion with the
Septuagint, in the Tetrapla'' (\work{Church History} 6.16; cf.\ Jerome,
\work{On Illustrious Men} 54). The purpose was controlled
comparison: to know, in argument with Jews and heretics, exactly what the
Church's Greek Bible shared with the Hebrew and where it differed. The
gigantic original---by any estimate thousands of pages---was almost
certainly never copied whole; it survived in the library of Caesarea, where
Jerome consulted it, and perished with that library, leaving fragments,
marginal collations, and the Syriac translation of one column.

\subsection{Scripture's senses, and the school's method}

On that textual base rose the exegesis, in three genres: scholia (brief
notes), homilies (preached exposition---several hundred survive, almost all
in Latin), and verse-by-verse commentaries, of which the greatest, on John
and Matthew, ran to dozens of books. Its governing conviction was that
Scripture is a single divine economy whose obvious sense opens onto deeper
ones---as body opens onto soul and spirit---so that the interpreter's task
is prayerful ascent, not mere philology, though never less than philology.
The program is stated in his own advice to a departing student: ``I wish to
ask you to extract from the philosophy of the Greeks what may serve as a
course of study or a preparation for Christianity,'' as Israel took the gold
of Egypt to build the sanctuary---``by spoiling the Egyptians,'' the things
of God's service are furnished; and then: ``Do you then, my son, diligently
apply yourself to the reading of the sacred Scriptures\ldots\ knock at its
locked door, and it will be opened to you'' (\work{Letter to Gregory}
1--3). Philosophy is preparation; Scripture, read on one's knees, is the
thing itself.

\subsection{\work{On First Principles}}

Before he left Alexandria Origen also wrote the most consequential and most
disputed of his works: ``He wrote also the books \work{De Principiis} before
leaving Alexandria'' (\work{Church History} 6.24.3)---\work{On First
Principles}, the first Christian attempt at a systematic investigation of
God, the rational creation, freedom, and Scripture. Its method is announced
in the preface and is essential to any fair judgment of its author. The
fixed points are what the Church teaches: ``that alone is to be accepted as
truth which differs in no respect from ecclesiastical and apostolic
tradition,'' the teaching ``transmitted in orderly succession from the
apostles, and remaining in the Churches to the present day.'' But the
apostles, Origen continues, deliberately left questions open---whether the
soul is transmitted by generation or comes to the body from without ``is not
distinguished with sufficient clearness in the teaching of the Church'';
of the devil and his angels ``what they are, or how they exist, it has not
explained with sufficient clearness''; ``what existed before this world, or
what will exist after it, has not become certainly known to the many''
(\work{On First Principles}, preface 2--7). Into exactly those open spaces
he sent his hypotheses: a first creation of rational spirits; their
pre-mundane cooling and fall, diversely deep, into the diversity of the
present order; the world as a just and remedial arrangement for their
recovery; and the hope---argued from Paul---of an end like the beginning.
``These subjects, indeed, are treated by us with great solicitude and
caution, in the manner rather of an investigation and discussion, than in
that of fixed and certain decision\ldots\ We think, indeed, that the
goodness of God, through His Christ, may recall all His creatures to one
end, even His enemies being conquered and subdued'' (\work{On First
Principles} 1.6.1). ``The end is always like the beginning'' (1.6.2): the
famous \textit{apokatastasis}, the restoration of all things, is proposed
there in precisely this exploratory register---which its author announced,
and which much of posterity ignored in both directions, hardening the
hypotheses into doctrines before condemning them.

\witnesslimit{\work{On First Principles} survives complete only in the Latin
of Rufinus of Aquileia (398), who states his method openly: a predecessor
translating Origen had ``so smoothed and corrected them in his translation,
that a Latin reader would meet with nothing which could appear discordant
with our belief,'' and Rufinus followed suit---``If, therefore, we have
found anywhere in his writings, any statement opposed to that view, which
elsewhere in his works he had himself piously laid down regarding the
Trinity, we have either omitted it, as being corrupt, and not the
composition of Origen, or we have brought it forward agreeably to the rule
which we frequently find affirmed by himself'' (Rufinus, preface to
\work{On First Principles}). Substantial Greek survives only in excerpts,
some transmitted by enemies. Every quotation of the work in this study is
therefore a quotation of the Rufinian layer, and no reconstruction of
Origen's exact Greek is claimed.}
